\originalchapter{8.333复习包3：关联气体与量子正则系综}
\osource{8.333 Fall 2013页面公开的Fall 2007 Test3 Review及Solutions，10题。题3、5对应已有卷内完整题解；同卷重复编号保留映射，不重复抄页。}
\begin{exercise}[原题1(a)--(g)]\label{r:333-3-1}
$N$电子及均匀正背景，$U_Q=\sum_{i<j}V(q_i-q_j)$、$V(r)=e^2/(4\pi r)-c$，背景令$\int V=0$。(a)Fourier非零模$e^2/q^2$、零模0；(b)证明$l$环积分$R_l=V^{1-l}\int_q\widetilde V(q)^l$；(c)累积量中环计数$S_l\simeq(l-1)!N^l/2$；(d)环求和；(e)积分得$\log Z-\log Z_0=V\kappa^3/(12\pi)$，$\kappa^2=\beta e^2N/V$；(f)压强；(g)两粒子间无环链的求和屏蔽势。
\end{exercise}
\begin{solution}
Coulomb的$-\nabla^2(1/4\pi r)=\delta$给$e^2/q^2$；背景专门去掉$q=0$，没有正负长程项残余。环卷积将各边动量固定为同一$q$，一个整体位置积分$V$，即(b)。选$l$不同顶点$\binom Nl$、无向环排列$(l-1)!/2$，累积量展开还含$l!$次边次序，结合给$N!/(N-l)!\,(l-1)!/2$；$l=2$重复同一对边直接验证$N(N-1)/2$。

$\sum_{l\ge2}(-x)^l/(2l)=[x-\log(1+x)]/2$，故
\begin{equation}
\Delta\log Z=\frac V2\int_q\left[\frac{\kappa^2}{q^2}-\log\left(1+\frac{\kappa^2}{q^2}\right)\right].
\end{equation}
把$\kappa$微分，$\partial_\kappa(\Delta\log Z/V)=\kappa\int_q[1/q^2-1/(q^2+\kappa^2)]=\kappa^2/(4\pi)$；在$\kappa=0$定积分常数0即$\kappa^3/(12\pi)$。固定$N,T$有$\kappa\propto V^{-1/2}$，压力修正$\Delta P=\kb T\partial_V\Delta\log Z=-\kb T\kappa^3/(24\pi)$，非解析$n^{3/2}$，不可用有限次数整数virial展开替代。

链边的几何级数$\widetilde V_{scr}=\widetilde V/(1+\beta n\widetilde V)=e^2/(q^2+\kappa^2)$，逆变换$V_{scr}=e^2\e^{-\kappa r}/(4\pi r)$。若题末写无$e^2$，那是势除以耦合后的归一化；本解保持原能量单位。环近似用弱Coulomb耦合，不声称所有密度下精确。
\end{solution}
\begin{exercise}[原题2(a)--(d)]\label{r:333-3-2}
$d$维阶梯势$V=\infty$（$r<a$）、$-\varepsilon$（$a<r<b$）、0（$r>b$）。(a)$B_2$及高低温；(b)压缩率首修正；(c)高温改写van der Waals式和参数；(d)$b=a,d=1$的$B_3$。
\end{exercise}
\begin{solution}
球体积$\Omega_dR^d$、$\Omega_d=\pi^{d/2}/\Gamma(d/2+1)$。$f=-1$于核、$\e^{\beta\varepsilon}-1$于壳，故$B_2=\Omega_d[a^d-(b^d-a^d)(\e^{\beta\varepsilon}-1)]/2$。高温趋核体积半值，低温吸引使它大负；温度低到关联团簇多，稀薄截断失控。

$P=n\kb T(1+B_2n+\cdots)$，$\kappa_T=[n\kb T(1+2B_2n)]^{-1}\simeq[1-2B_2n]/(n\kb T)$。高温$B_2=b_{vdW}-a_{vdW}/(\kb T)$，$b_{vdW}=\Omega_da^d/2$、$a_{vdW}=\Omega_d\varepsilon(b^d-a^d)/2$，重求和$P\simeq n\kb T/(1-b_{vdW}n)-a_{vdW}n^2$仅匹配到$n^2$。一维硬棒精确$P=n\kb T/(1-na)$，$B_3=a^2$，见习题\ref{p:334-4-4}。
\end{solution}
\originalsection{原题3：Dieterici临界点的完整引用}
原题$P(v-b)=\kb T\exp[-a/(\kb Tv)]$三问：临界压缩因子、$v_c$上压缩率、临界等温线最低非零项，恰为习题\ref{p:334-1-1}(e)--(g)。所有条件及解答均在该处，复习编号保留为333-R3-3。
\begin{exercise}[原题4(a)--(e)]\label{r:333-3-4}
二维$N$个$+c_0$、$N$个$-c_0$经典粒子，面积$A=L^2$，$H=\sum_{1}^{2N}p_i^2/(2m)-\sum_{i<j}c_ic_j\log(|q_i-q_j|/r_0)$。(a)吸、斥对数量；(b)$Z$积分并化成距离幂；(c)用$q'=q/L$证明$Z\propto A^{2N-\beta Nc_0^2/2}$；(d)二维压力与低温问题；(e)加入硬核$a$，由$N=1$估其重要温度与相的性质。
\end{exercise}
\begin{solution}
同号斥对$2\binom N2=N(N-1)$，异号吸对$N^2$。动量给$\lambda^{-4N}$，$Z=[(N!)^2\lambda^{4N}]^{-1}\int_A\prod\dd^2q_i\prod_{i<j}(|q_i-q_j|/r_0)^{\beta c_ic_j}$，固定$r_0$只是指定能零。总对积电荷$\sum_{i<j}c_ic_j=[(\sum c)^2-\sum c_i^2]/2=-Nc_0^2$。坐标缩放使测度$L^{4N}$、距离幂$L^{-\beta Nc_0^2}$，即题示面积幂，前提积分收敛。

因此$P=N(2\kb T-c_0^2/2)/A$，高温回到$2N$粒子理想压，形式$\kb T<c_0^2/4$变负。但没有核的异号吸引在更早的短距条件$\beta c_0^2\ge2$便使两体积分$\int_0^Lr^{1-\beta c_0^2}\dd r$发散，形式负压分支并不存在稳定配分函数。

硬核下$\int_a^Lr^{1-\beta c_0^2}\dd r$在$\beta c_0^2<2$由大距控制、$>2$由核控制，边界对数。所以单对估计温标$\kb T_*=c_0^2/2$，低温偏好束缚中性对、高温自由电荷。有限密度的集体涡旋/Coulomb气转变需屏蔽与逸度RG，单对核重要尺度不等于已经精确证明多体BKT温度。
\end{solution}
\originalsection{原题5：最近邻一维气的完整引用}
排序粒子硬核最近邻势，原题(a)--(f)逐项对应习题\ref{p:334-4-4}：Hamilton量、间距域、Laplace鞍点、边界间隙积分、平均长度、硬棒方程与$B_l$。原复习题另问排除体积比较：一维$b_{vdW}=a$与精确硬棒结果相合；有限维一般van der Waals排除重求和并非精确。
\begin{exercise}[原题6(a)--(d)]\label{r:333-3-6}
$N+1$质量$m$粒子以$N$弹簧$K$、原长$a$连成链，端点固定相距$Na$，位移$u_0=u_N=0$，$H=\sum_{1}^{N-1}p_i^2/(2m)+K\sum_{i=0}^{N-1}(u_{i+1}-u_i)^2/2$。(a)正弦正规模与频率；(b)经典$Z$；(c)$\avg{|\widetilde u_k|^2},\avg{u_i^2}$并画位置曲线；(d)仅左端固定、右端自由如何改变。
\end{exercise}
\begin{solution}
$u_i=\sqrt{2/N}\sum_{r=1}^{N-1}\widetilde u_r\sin(\pi ri/N)$是实正交变换；频率$\omega_r=2\sqrt{K/m}\sin(\pi r/(2N))$。独立模$H_r=\widetilde p_r^2/(2m)+m\omega_r^2\widetilde u_r^2/2$，$Z=\prod_r\kb T/(\hbar\omega_r)=[\kb T/(\hbar\sqrt{K/m})]^{N-1}/\sqrt N$，乘积恒等式$\prod_{r=1}^{N-1}2\sin(\pi r/(2N))=\sqrt N$。

$\avg{\widetilde u_r^2}=\kb T/(m\omega_r^2)$；离散Laplace逆矩阵给$\avg{u_iu_j}=(\kb T/K)[\min(i,j)-ij/N]$，所以$\avg{u_i^2}=(\kb T/K)i(N-i)/N$，中部最大，随链长线性增长。
\begin{center}\begin{tikzpicture}\begin{axis}[width=9cm,height=4cm,xlabel={$i/N$},ylabel={$K\avg{u_i^2}/(N\kb T)$},xmin=0,xmax=1,ymin=0]\addplot[thick,structurecolor,domain=0:1]{x*(1-x)};\end{axis}\end{tikzpicture}\end{center}
自由右端保留$u_1,\ldots,u_N$，独立弹簧伸长$\delta_i=u_i-u_{i-1}$，$\avg{\delta_i\delta_j}=\kb T\delta_{ij}/K$，故$\avg{u_iu_j}=\kb T\min(i,j)/K$、$Z=[\kb T/(\hbar\sqrt{K/m})]^N$。频率$2\sqrt{K/m}\sin[(2r-1)\pi/(4N+2)]$，$r=1,\ldots,N$。原提示说$N-1$伸长是沿用固定端自由度的索引；若终点真自由，共有$N$个移动粒子与伸长，不能漏最后一个模。
\end{solution}
\begin{exercise}[原题7(a)--(g)]\label{r:333-3-7}
黑洞熵$S=\kb c^3A/(4G\hbar)$。(a)牛顿逃逸速设为$c$得半径；(b)两个太阳质量$M_\odot\simeq2\times10^{30}$ kg黑洞合并熵增、比特数；(c)$E=Mc^2$给温度；(d)黑体辐射功；(e)蒸发时长和太阳质量估值；(f)与题给2.7 K背景平衡质量；(g)半径$R$球内全息熵上界。采用题设无自旋黑体理想模型，不加未给灰体和额外粒子种类。
\end{exercise}
\begin{solution}
$v_{esc}^2=2GM/R$，$R_s=2GM/c^2$，$A=16\pi G^2M^2/c^4$，所以$S=4\pi\kb GM^2/(\hbar c)$。忽略合并辐射，末质量$2M$，熵增$8\pi\kb GM^2/(\hbar c)$，约$3.1\times10^{77}$ bits（除$\kb\log2$）；真实计算应包括辐射能与熵，不能把无损质量相加作为精确天体结果。

$1/T=\partial S/\partial E$给$T_H=\hbar c^3/(8\pi\kb GM)$，太阳质量约$6.1\times10^{-8}$ K。用黑体$P=\sigma_{SB}AT_H^4=\hbar c^6/(15360\pi G^2M^2)$，$\dot Mc^2=-P$积分得到$t_{ev}=5120\pi G^2M_0^3/(\hbar c^4)$，太阳质量约$2.2\times10^{67}$年。末期黑体半经典模型失效，此巨大尺度只作原题估算。

$T_H=2.7$ K给$M\simeq4.5\times10^{22}$ kg。半径$R$所容最大熵按面积界为$S_{max}=\pi\kb c^3R^2/(G\hbar)$，体现面积而非普通体积增长；界的引力和几何前提属于题设全息原理，不由经典相空间计数单独证明。
\end{solution}
\begin{exercise}[原题8(a)--(h)]\label{r:333-3-8}
量子谐振子$H=p^2/(2m)+m\omega^2q^2/2$。(a)$Z,U$；(b)本征基密度矩阵；(c)非对易算符指数导数与迹导数；(d)用$\partial_mZ=0$证动势能相等；(e)$\avg{q^2}$及链低温量子修正；(f)高温坐标核；(g)低温核；(h)完整精确坐标核。
\end{exercise}
\begin{solution}
$b=\beta\hbar\omega$，$Z=\sum_{n\ge0}\e^{-b(n+1/2)}=[2\sinh(b/2)]^{-1}$，$U=(\hbar\omega/2)\coth(b/2)$，$\rho=Z^{-1}\sum_n\e^{-\beta E_n}\ket n\bra n$。

任意$A(x)$有$\partial_x\e^A=\int_0^1\e^{(1-s)A}A'\e^{sA}\dd s$，由指数乘积展开或Duhamel式导出，只有$[A,A']=0$才可写$A'\e^A$。迹内循环却始终给$\partial_x\Tr\e^A=\Tr(A'\e^A)$。固定$\omega$时$Z$与$m$无关，所以$0=\avg{\partial_mH}=\avg{-p^2/(2m^2)+\omega^2q^2/2}$，两能各$U/2$，$\avg{q^2}=\hbar\coth(b/2)/(2m\omega)$。

链的每模将$\kb T/(m\omega_r^2)$换为$\hbar\coth(\beta\hbar\omega_r/2)/(2m\omega_r)$，再按正弦模相加；零温仍有零点涨落，一维长波$1/|k|$给对数链长而非经典线性增长。

完整核可从热核方程推导：设$\langle q|\e^{-\beta H}|q^\prime\rangle$的非对角一般形式为$A_\beta\exp[-a_\beta(q^2+q^{\prime2})+b_\beta qq^\prime]$。代$-\partial_\beta K=H_qK$并比较多项式系数，得$a_\beta'=m\omega^2/2-2\hbar^2a_\beta^2/m$、$b_\beta'=-2\hbar^2a_\beta b_\beta/m$、$A_\beta'/A_\beta=-\hbar^2a_\beta/m$。$\beta\to0$的自由Gaussian $\delta$核选出$a=m\omega\coth b/(2\hbar)$、$b_\beta=m\omega/(\hbar\sinh b)$、$A=\sqrt{m\omega/(2\pi\hbar\sinh b)}$。再除$Z$、重写中点和差坐标，给
\begin{equation}
\rho(q',q)=\sqrt{\frac{m\omega}{\pi\hbar}\tanh\frac b2}\exp\left[-\frac{m\omega}{\hbar}\tanh\frac b2\left(\frac{q+q'}2\right)^2-\frac{m\omega}{4\hbar}\coth\frac b2(q-q')^2\right].
\end{equation}
高温$b\ll1$是$\sqrt{\beta m\omega^2/(2\pi)}\exp[-\beta m\omega^2(q+q')^2/8-m(q-q')^2/(2\beta\hbar^2)]$到领先阶；低温$b\gg1$是$\psi_0(q')\psi_0(q)$，$\psi_0=(m\omega/\pi\hbar)^{1/4}\e^{-m\omega q^2/(2\hbar)}$。沿对角积分1、对角二阶矩等于(e)，为独立归一化检查。
\end{solution}
\begin{exercise}[原题9(a)--(f)]\label{r:333-3-9}
三维$N$正、$N$负相对论量子粒子$H=c\sum_i|p_i|+\sum_{i<j}e_ie_j/|r_i-r_j|$，立方边$L$。(a)Schrödinger式与Bose/Fermi对称；(b)$E_n(L)$缩放；(c)$Z$的$T,V$组合；(d)$E,P$精确关系；(e)$d$维Coulomb$r^{2-d}$配非相对论$p^2/(2m)$何维同样齐次及关系；(f)为何密集真实Coulomb混合物可能不服这些式。
\end{exercise}
\begin{solution}
$[c\sum_i\sqrt{-\hbar^2\nabla_i^2}+\sum e_ie_j/r_{ij}]\Psi_n=E_n\Psi_n$，同正电荷组、同负组分别按全粒子含自旋的Bose对称或Fermi反对称，不对不同电荷强行交换。$r_i=Lr_i'$、规范化$\Psi_L=L^{-3N}\Psi_1$使每动能、势能均$L^{-1}$，故$E_n(L)=E_n(1)/L$，前提Hamilton量与边界确无另一长度且稳定谱存在。

$Z=\sum_n\exp[-E_n(1)/(L\kb T)]$仅依赖$LT$，$E=-\partial_\beta\log Z$，$P=\kb T\partial_V\log Z=E/(3V)$。非相对论动能$L^{-2}$，须$d-2=2$即$d=4$势也$L^{-2}$，于是$E=dPV/2=2PV$。

\textbf{同一个稳定齐次Hamilton模型}不会仅因变液体或固体就丧失这个形式恒等式。真实密集材料的核尺寸、电子质量/非相对论转换、排斥核、有效相互作用与环境引入额外尺度；异号点吸引也可能导致塌缩、无有限$Z$。是这些前提变化限制缩放，不是“高密度”三个字自动使数学证明失效。
\end{solution}
\begin{exercise}[原题10(a)--(d)]\label{r:333-3-10}
一般$H(p,q)$。(a)经典正则配分在一对$q_1\to\lambda q_1,p_1\to p_1/\lambda$的canonical变换下不变；(b)量子对应；(c)$H=\sum p_i^2/(2m)+V(q)$证明$\avg{p_1^2/m}=\avg{q_1\cdot\nabla_1V}$；(d)虚构G-8.333星系边缘速200 km/s，估质量/尺度比。
\end{exercise}
\begin{solution}
相空间全空间或连同边界一并变换时Jacobian$\lambda^d\lambda^{-d}=1$，$Z[H_\lambda]=Z[H]$。量子用幺正dilation$U_\lambda$保持$[q,p]=\ii\hbar$，$H_\lambda=U_\lambda H U_\lambda^{-1}$谱不变，$\Tr\e^{-\beta H_\lambda}$不变。不能在固定硬盒里缩坐标却忽略域/墙项；这时边界压力会贡献virial。

在$\lambda=1$微分$\log Z$，迹Duhamel式适用于量子非对易，$0=-\beta\avg{-p_1^2/m+q_1\cdot\nabla_1V}$，即题式。星系$Mv^2\sim GM^2/R$给$M/R\sim v^2/G\simeq6.0\times10^{20}$ kg/m，约$9.3\times10^9M_\odot$每kpc。质量分布、轨道各向异性与边缘速定义带来常数阶差，题给速度只能估尺度，不确定完整质量剖面。
\end{solution}
